\chapter{Paired-Spectrum Quotient Correspondence}
\label{chap:paired-spectrum-quotient-correspondence}

\status{Exact correspondence theorem proved. The finite paired xi-zero set and finite paired $F$-root set are related by an exact two-to-one quotient. SOH-G003 and RH remain open.}

Chapters~\ref{chap:negative-inversion-zero-set-no-go} and
\ref{chap:quotient-negative-inversion-zero-set-no-go} establish two finite
exceptional paired spectra: one in the $s$-plane and one in the even quotient
$w$-plane.  This chapter proves that they are not merely analogous finite
sets.  They are the same exceptional structure before and after the exact
two-sheeted functional-reflection quotient.

\section{The even quotient map}

Define
\begin{equation}
 q(s)=\left(s-\frac12\right)^2.
 \label{eq:g016-q}
\end{equation}
Then
\begin{equation}
 q(1-s)=q(s),
 \label{eq:g016-reflection-quotient}
\end{equation}
so $q$ identifies precisely the pair exchanged by the holomorphic functional
reflection $s\mapsto1-s$.

For $w\ne0$, the fiber consists of the two distinct points
\begin{equation}
 q^{-1}(w)
 =\left\{\frac12+\sqrt w,\frac12-\sqrt w\right\}.
 \label{eq:g016-fiber}
\end{equation}
The two square-root choices are exchanged by $s\mapsto1-s$.

From Chapter~\ref{chap:positive-riemann-kernel},
\begin{equation}
 F(0)=\xi(1/2)>0,
 \label{eq:g016-zero-not-root}
\end{equation}
so $0\notin Z_F$.  Consequently every quotient root relevant below has
exactly two distinct preimages.

\section{The commuting negative-inversion diagram}

G014 gives
\[
 N_s(s)=\frac{s-1}{2s-1},
\]
while G015 gives
\[
 J(w)=\frac1{16w}.
\]
Writing $z=s-1/2$, G012 gives
\[
 N_z(z)=-\frac1{4z}.
\]
Squaring immediately yields
\begin{equation}
 \boxed{
 q(N_s(s))=J(q(s)).
 }
 \label{eq:g016-commuting-diagram}
\end{equation}
Thus
\[
 \boxed{q\circ N_s=J\circ q.}
\]
The quotient and negative inversion commute exactly; there is no fitted or
asymptotic step in this diagram.

\section{The two paired spectra}

Recall the G014 paired xi-zero set
\begin{equation}
 P_N
 =\{\rho:\xi(\rho)=0,\ \xi(N_s(\rho))=0\}
 \label{eq:g016-PN}
\end{equation}
and the G015 paired quotient-root set
\begin{equation}
 P_J
 =\{w:F(w)=0,\ F(J(w))=0\}.
 \label{eq:g016-PJ}
\end{equation}
Both sets are finite by the preceding no-go theorems.

\begin{theorem}[Exact paired-spectrum quotient]
\label{thm:g016-paired-quotient}
The quotient map satisfies
\begin{equation}
 \boxed{q(P_N)=P_J.}
 \label{eq:g016-image-equality}
\end{equation}
Moreover,
\[
 q:P_N\longrightarrow P_J
\]
is exactly two-to-one.  Hence
\begin{equation}
 \boxed{|P_N|=2|P_J|.}
 \label{eq:g016-cardinality}
\end{equation}
\end{theorem}

\begin{proof}
Take $\rho\in P_N$.  Since $\xi(\rho)=0$, the even quotient identity
\[
 \xi\!\left(\frac12+z\right)=F(z^2)
\]
gives
\[
 F(q(\rho))=0.
\]
Also $\xi(N_s(\rho))=0$, so by
\cref{eq:g016-commuting-diagram},
\[
 F(J(q(\rho)))
 =F(q(N_s(\rho)))
 =0.
\]
Therefore $q(\rho)\in P_J$, proving $q(P_N)\subseteq P_J$.

Conversely, let $w\in P_J$.  Equation~\eqref{eq:g016-zero-not-root} ensures
$w\ne0$.  The two points in \cref{eq:g016-fiber} both satisfy
\[
 \xi(\rho_\pm)=F(w)=0.
\]
By the commuting diagram,
\[
 q(N_s(\rho_\pm))=J(w).
\]
Since $F(J(w))=0$, every point in the two-element fiber over $J(w)$ is an xi
zero.  In particular,
\[
 \xi(N_s(\rho_\pm))=0.
\]
Thus both $\rho_+$ and $\rho_-$ lie in $P_N$, and $w\in q(P_N)$.  Hence
$P_J\subseteq q(P_N)$ and equality follows.

Finally, every $w\in P_J$ is non-zero, so its fiber consists of exactly two
distinct points, and the preceding argument places both in $P_N$.  Therefore
the restriction of $q$ is exactly two-to-one, proving
\cref{eq:g016-cardinality}.
\end{proof}

\section{Lift of the exceptional orbit structure}

Both $N_s$ and $J$ are involutions.  Equation~\eqref{eq:g016-commuting-diagram}
therefore lifts every $J$-orbit in $P_J$ to $N_s$-orbits in $P_N$.

If $\{w,J(w)\}$ is a non-fixed two-cycle, its two quotient fibers contain four
xi zeros.  Negative inversion arranges these four points into two
$N_s$ two-cycles.  The functional reflection exchanges the two lifted cycles.

The only possible fixed quotient root left by G015 is $w=-1/4$.  Its fiber is
\begin{equation}
 q^{-1}(-1/4)
 =\left\{\frac12+\frac{i}{2},\frac12-\frac{i}{2}\right\},
 \label{eq:g016-minus-quarter-fiber}
\end{equation}
and G012 proves that both points are individually fixed by $N_s$.
The positive fixed quotient value $+1/4$ contributes no paired fiber because
$F(1/4)>0$.

\section{Cardinality structure}

Let $M$ denote the number of non-fixed two-cycles in $P_J$, and let
\[
 \varepsilon\in\{0,1\}
\]
record whether $F(-1/4)=0$.  Then the G015 orbit classification and
\cref{thm:g016-paired-quotient} give
\begin{equation}
 |P_J|=2M+\varepsilon,
 \qquad
 |P_N|=4M+2\varepsilon.
 \label{eq:g016-cardinality-structure}
\end{equation}
In particular,
\begin{equation}
 \boxed{|P_N|\text{ is even}.}
 \label{eq:g016-PN-even}
\end{equation}
The unresolved status of $F(-1/4)$ changes only the final $+2$ term in the
xi-layer paired cardinality.  It does not affect the exact two-to-one
correspondence.

\section{Proof firewall}

Proved in SOH-G016:
\begin{itemize}
 \item $q(s)=(s-1/2)^2$ identifies the functional-reflection pair;
 \item $q\circ N_s=J\circ q$ exactly;
 \item $q(P_N)=P_J$;
 \item $q:P_N\to P_J$ is exactly two-to-one;
 \item $|P_N|=2|P_J|$ and therefore $|P_N|$ is even;
 \item every non-fixed paired quotient two-cycle lifts to two paired xi-zero
       two-cycles;
 \item the possible $w=-1/4$ fixed root lifts to the G012 fixed xi pair.
\end{itemize}

Not proved or claimed in SOH-G016:
\begin{itemize}
 \item whether $F(-1/4)$ vanishes;
 \item that either paired set is empty;
 \item real-rootedness of $F$;
 \item PF$_\infty$;
 \item SOH-G003;
 \item the Riemann Hypothesis.
\end{itemize}
